\section{What Would Show This Is Wrong}
\label{sec:study}

Three claims in this paper fail independently of one another, and each one fails
in a way that a different kind of evidence catches. The record either reproduces
the bytes the developer's editor and agent left in the files or it does not, and
that is a question about a program, settled by running it over histories that
already exist. The grouping either lines up with the requests a developer made or
it does not, and that is a question about an inference, settled by comparing it
against sessions where somebody happened to write the requests down. A developer
either gets more of these operations right with \sgt{} than without it, or does
not, and that is a question about people. The first has a procedure and a number we
can print today, the second is under way on a corpus we are still assembling, and
the third is designed and not yet run. Underneath all three, \sgt{} has been the
version control for its own repository since the day it first ran, which is where
every number in Section~\ref{sec:discussion} comes from and why the two prices
Section~\ref{sec:design} states most flatly are the two that cost us something
before they cost anybody else.

\smallskip\noindent\emph{Not one number in this section comes from that third
study.} The measure that decides the argument is not the time a task takes but the
errors a participant does not notice, and naming that before the data exists is the
only way to be believed about it afterwards.

\subsection{The rebuild has to be exact}
\label{sec:eval-exact}

If \sgt{} reassembles a file into bytes that differ from what the developer left
there, then every operation in Section~\ref{sec:design} is set arithmetic over a
fiction, and nothing else in this paper is worth reading. This is the one
measurement whose failure would end the work rather than qualify it, so it goes
first and it is the only one we insist on before adoption.

The procedure follows the shape of the claim. We walk a repository's history one
commit at a time, mine each commit into records, rebuild every file from the set of
records that commit's state implies, and compare the result against the blob git
holds, byte for byte. A repository passes at a commit only if every file matches, so a
single differing byte in a single file is a failure at that commit, and we will report
the fraction of files that rebuild rather than an average over repositories, because a
developer cares whether their file comes back and not whether most files do. The
corpus is open source Python and TypeScript repositories chosen for length of history
and nothing else, and we are still assembling it, so the only number this procedure has
produced comes from our own repository, which is the hardest case available to us
because records written by five different versions of our miner sit in it.

That number is fifty-six files that no valid set of records reproduces, one in seven of
the files whose records name functions, and Section~\ref{sec:unreproducible} gives that
denominator and says what is in the fifty-six. What we do not know is whether
fifty-six is the cost of our own format changes or the cost of the design, and since
five miner versions in fifty-seven days is a rate no adopter will ever match, the
corpus run is how we intend to separate the two: a repository \sgt{} has never touched
cannot have our migrations in it.

\subsection{Whether the grouping matches what was asked for}

The names in every figure in this paper are inferred, and a developer who types
\cmd{sgt revert "price cache"} is trusting that inference with the contents of
their files. The question is not whether the clustering scores well against a
partition somebody drew, which is how the untangling literature is evaluated for
the good reason that its authors had nothing better to compare against
(Section~\ref{sec:related}). The question is whether the set that comes back when
a developer names one request is the set of edits that request produced.

That question needs a corpus where the requests were recorded, which agent sessions
now supply as a matter of course, since the transcript of a session names each request
and the files the agent touched while carrying it out. We mine each repository without
reading any transcript, derive a reference grouping from the transcripts alone, and
then ask two things of every request in the session: how much of what that request
produced comes back when the developer names it, and how much else comes with it.
Those are the two ways a revert disappoints the person who ran it, the edit it missed
and left sitting in the file, and the edit belonging to a different request that it
carried off as well. They fail asymmetrically, and the second is worse, because a
developer who is handed too little notices when the code still runs the way it did and
a developer who is handed too much has already lost work they meant to keep.

We expect the shared helper to be where this measurement goes wrong, because
Section~\ref{sec:names} predicts it and fifty-seven days of daily use have not
stopped it happening. An argument parser or a logging function that changes in nearly every
save is coupled to everything by the signal the clustering reads, and it draws
otherwise unrelated requests into one group. If most of the disagreement between
our grouping and the transcripts turns out to sit on functions of that kind, then
the fix is to weight a function's coupling down as the number of things it couples
to grows, and that is a change to one term rather than to the design. If the
disagreement is spread evenly across ordinary functions instead, then the coupling
signal is not carrying the grouping and the developer should be asked at save
time, which is the design we argued against in Section~\ref{sec:names} and would
have to argue for.

\subsection{Three requests, a transcript, and code nobody in the room wrote}
\label{sec:eval-study}

\begin{table*}[t]
  \centering
  \caption{The six measures of the study in Section~\ref{sec:eval-study}, each
  with the result that would count against the design. The first row is the one we
  would report first, and the third row is the one a reader would expect us to
  report first.}
  \Description{A three-column table. The left column names six measures: undetected
  errors, correct submissions, time to completion, diff lines read, forks raised and
  judged pointless, and label corrections. The middle column states what the authors
  expect for each. The right column states, for each, the observation that would
  count against the design.}
  \label{tab:measures}
  \small
  \begin{tabular}{@{}p{0.235\textwidth}p{0.335\textwidth}p{0.355\textwidth}@{}}
    \toprule
    \textbf{What we measure} & \textbf{What we expect} &
      \textbf{What would count against us} \\
    \midrule
    Submissions the participant believed correct and that were not.
      & The gap here is wider than the gap in time. Removing a request touches
        code the participant never read, and in git nothing tells them so.
      & Equal rates in both conditions. The recovery we call expensive would then
        be something developers already do reliably, and Section~\ref{sec:scenarios}
        is a story about us. \\[0.6ex]
    Submissions that pass the check.
      & More of them with \sgt{}, and the ones that fail fail in the places
        \cmd{revert} named out loud.
      & A failure in \sgt{} that the preview did not name. A report that misses a
        consequence is worse than no report. \\[0.6ex]
    Time to the moment the participant says they are done.
      & Shorter with \sgt{}, and we would not publish this alone.
      & Shorter with \sgt{} while the first row is flat. \sgt{} would then be a
        convenience, and we have argued for it as something else. \\[0.6ex]
    Lines of diff the participant opened before submitting.
      & Far fewer with \sgt{}, since no command in
        Section~\ref{sec:design} asks anyone to look at a hunk.
      & Participants reading the diff anyway. They would be checking \sgt{}'s
        report rather than using it, and after twenty minutes of that they are
        doing the git task with extra steps. \\[0.6ex]
    Forks \sgt{} raised, and how many the participant judged pointless.
      & A small number, mostly judged pointless, all cheap.
      & Any participant switching the rule off, or a fork resolved by picking a
        side whose result then failed the check. Being stopped would have bought
        nothing. \\[0.6ex]
    Labels the participant corrected, and what the correction cost.
      & One or two renames, no regroups.
      & A regroup on the critical path of an operation. The inference would then
        be spending the developer's attention rather than saving it. \\
    \bottomrule
  \end{tabular}
\end{table*}

The study puts a developer in the position the paper is about, which means giving
them a repository whose code they did not write and the three sentences that
caused it. Each participant receives a five thousand line Python service, a
transcript of the session that produced its most recent hour of work, and three
tasks, each of which names one request from that transcript and asks for it to be
moved or removed without disturbing the other two. This is the
state of the developer in Section~\ref{sec:scenarios} a week later, and it is also
the state of the developer who picks up a colleague's branch on a Monday, and it
is the only arrangement we found in which every participant knows exactly as much
about the code as every other one. Participants who direct the agent themselves
each remember a different amount of what came back, and that remembering is the
thing whose absence we are claiming matters.

We recruit developers who use a coding agent as part of their ordinary work,
because the difficulty we describe belongs to them and to nobody else, and we ask
each of them to work through two repositories, one with git and one with \sgt{},
in a counterbalanced order. In the git condition they use whatever they normally
use, including a graphical client, since the argument is about what git's record
can express and not about anybody's command line fluency. Each repository holds
three requests whose edits interleave inside shared functions, and each has a
different set of three so that a participant who has learned the shape of the
first task cannot apply it to the second. The tasks are the three operations from
Section~\ref{sec:scenarios} with a checkable outcome: take one named request out
and leave the other two working, carry one request to a release branch on its own,
and find the earlier state of a request in order to change it. We leave the other
two scenarios out of the study because their outcome is a judgement about when to
commit and how to stack, and we would be scoring participants against our
preference rather than against a result.

A submission is correct when a script says so, and the script asks three
questions. The tests belonging to the removed request fail, the tests belonging to
the two kept requests pass, and no function outside the intended set differs from
the reference repository. That third question is the one that catches the failure
this paper is about, because a participant who removes the caching and quietly
takes half the retry policy with it has produced a repository that passes its own
tests.

The first row of Table~\ref{tab:measures} is the one that matters, because a
developer who spends an extra twenty minutes untangling a commit by hand has lost
twenty minutes, while a developer who submits a repository they believe is clean and
that still calls a function they deleted has lost something they will not find until
the release branch fails to import. Hand untangling fails silently, and a design whose
whole argument is that the developer no longer holds the decomposition has to be
measured against silent failure rather than against the clock. We collect time because
a reader would want it and because a tool that is correct and unbearably slow is not
adopted, and we would not publish a time difference as the finding.

The last two rows of Table~\ref{tab:measures} measure the prices we stated in
Section~\ref{sec:design} rather than the benefits, and they are there because a
study that measures only what a design buys is a study that cannot be wrong. A
fork stops the developer and a wrong label makes them type a rename, both of which
are countable inside the session, and both have a threshold past which we would
change the design rather than defend it.

\subsection{What none of this would tell us}

Only the first of the three claims can be pushed past by running more of the same,
since adding repositories to Section~\ref{sec:eval-exact} sharpens that number and
nothing else here improves with volume.

A recorded request is not an intention. A developer who types
\emph{make the price lookup faster} holds a view of what should change that the
sentence does not contain, and the agent's response is a guess at the remainder.
The alignment we measure against a transcript is therefore alignment with what was
said, and the transcript sets a ceiling on it that no better clustering can pass.
We think this is the right measurement anyway, on the ground that what was said is
what the developer will type when they come back to remove it, but a reader who
believes intention and utterance come apart more often than we do should read that
number as an upper bound on something smaller.

Two hours with a transcript is not a week of forgetting. The cost we describe in
Section~\ref{sec:scenarios} grows as a developer's memory of the session fades,
and our design substitutes reading somebody else's transcript for that fading.
A participant who has just read three sentences knows more about the session than
the developer who typed them a week ago, so the git condition in our study is
easier than the situation it stands for, and the difference we measure understates
the one that matters. We prefer the understatement to a study nobody can run.

Finally, every task in the study has an answer a script can check, and we chose
those three tasks for that reason. The operations a developer wants most may be
the ones with no checkable answer, and deciding whether a stack of two reviews
was the right shape is one of them. We can say what \sgt{} makes possible there
and we cannot measure it, and Section~\ref{sec:discussion} keeps that distinction.
